\documentclass[12pt]{article}
\usepackage{fontspec}
\usepackage{polyglossia}
\setmainlanguage{french}
\usepackage{csquotes}

\usepackage{tikz}% faire des schémas

\begin{document}

\section{Un \emph{stemma codicum}}

% Chaque child {node {...}} est un enfant du nœud qui le précède.
% 2. D est un troisième enfant de A : on ajoute un child au même niveau que B et C.
% 3. E est un enfant de D : on place un child à l'intérieur des accolades de D.
% L'arbre entier est une seule instruction : un seul point-virgule, tout à la fin.
% 4. Le tikzpicture est placé dans un flottant figure, centré, avec une légende.

\begin{figure}[ht]
	\centering
	\begin{tikzpicture}
		\node {A}
			child {node {B}}
			child {node {C}}
			child {node {D}
				child {node {E}}
			}
		;
	\end{tikzpicture}
	\caption{Un \emph{stemma codicum}}
\end{figure}

\end{document}
